% ============================================================================================
\section{Introduction}
\label{sec:intro}
Assignment~2 asks each team to practise scientific visualization (SciVis) on a gridded geophysical dataset and
information visualization (InfoVis) on the tabular dataset it used in Assignment~1 (A1). There are three tasks per
domain. Each member owns exactly one SciVis and one InfoVis task, with no joint ownership apart from the shared data
preparation (Table~\ref{tab:contrib}).

\textbf{SciVis.} We use gridMET~\cite{abatzoglou2013,gridmet_site} for the three-month slice 1~July--30~September 2024
(92 days). % TODO verify: slice confirmed via the TA Google form
Late summer 2024 was meteorologically rich over CONUS. Four landfalling Atlantic hurricanes (Beryl, Debby, Francine and
Helene~\cite{nhc_tcr}) and a long western heat wave give spatially sharp, physically interpretable structure in all
three quantities we visualize: temperature (S1), atmospheric dryness (S2) and wind (S3). A quiet summer would be dominated
by smooth climatological gradients. These events instead put local extremes next to large-scale patterns, which is
exactly what tests colour normalisation, contour placement and vector sampling.

\textbf{InfoVis.} Our A1 dataset holds the roughly 1,000 all-time top posts of each of 50 large subreddits, 49,266 posts in
total, posted 2011--2024~\cite{kaggle_reddit}. Each post has a score (net upvotes), upvote ratio, comment and crosspost
counts, a timestamp, a post type and a link domain. None of the required structures (a network, hierarchies, a
multivariate table and time series) exists as such in the data. Section~\ref{sec:data} documents how each was derived.

\textbf{Tooling.} The work uses Python with xarray~\cite{hoyer2017}, Matplotlib~\cite{hunter2007}, Cartopy~\cite{cartopy},
scikit-image~\cite{vanderwalt2014} and NetworkX~\cite{hagberg2008}, plus D3.js v7~\cite{bostock2011} for the
browser views and Playwright~\cite{playwright} for scripted high-resolution screenshots. Every figure regenerates with one
command per task (README). Figures are numbered as in \fpath{images/}, and the figure index is
Table~\ref{tab:figindex} in the appendix.

\begin{table}[!t]
\centering
\caption{Contribution table. Each member owns one SciVis and one InfoVis task; only data preparation is shared.}
\label{tab:contrib}
\footnotesize
\begin{tabular}{p{1.9cm}p{2.9cm}p{2.8cm}}
\toprule
Member & Owned tasks & Shared data-preparation part \\
\midrule
Sidhanth Prabhu (BT2024027) & S1 colour mapping (+ GIFs); I1 node-link diagram, layouts, filtering, GEXF/PyVis & gridMET download, schema inspection, subsetting and unit conversion \\
Shrey Modi (BT2024125) & S2 Marching Squares, filled contours (+ GIF); I3 horizon and ridgeline plots & A1 cleaning; network, hierarchy, multivariate and time-series derivations \\
Parthsarati Samanta (BT2024083) & S3 quiver and streamlines (+ GIF); I2 D3 parallel coordinates and zoomable treemap & Wind decomposition, statistics tables and evidence-based date selection \\
\bottomrule
\end{tabular}
\end{table}
% TODO verify: split of the shared preprocessing between members (assigned by us; edit if different)

\subsection{The gridMET dataset}
gridMET combines the spatial detail of PRISM climatologies with the daily temporal attributes of the North American Land
Data Assimilation System (NLDAS-2)~\cite{mitchell2004nldas} to give daily surface meteorology on a
$1/24^\circ$ (about 4\,km) grid~\cite{abatzoglou2013}. Our files cover 25.07--49.40\,$^\circ$N and 124.77--67.06\,$^\circ$W
with $585\times1386$ cells, of which 483,879 (59.7\,\%) are land cells inside the CONUS mask. Wind speed and direction in gridMET
come from NLDAS-2 at about 1/8$^\circ$ and are interpolated to the 4\,km grid,
so wind is smoother than temperature or precipitation. Every value is a \emph{daily} aggregate: a daily mean
for wind and VPD, a daily maximum for temperature and humidity, and a daily total for precipitation. The global attributes state that a gridMET day
ends at midnight MST (07~UTC of the next calendar day). This matters for timing a hurricane landfall
(Sec.~\ref{sec:dates}).

\subsection{Choice of the eight time instances}
\label{sec:dates}
The assignment allows uniform or random sampling of 5--10 instances. We use a \emph{hybrid} rule so that the dates are
both spread over the season and centred on interpretable events. Each candidate event was first checked in the
data, not assumed. For every day we computed (Fig.~\ref{fig:01}, \fpath{daily_evidence.csv}):
\begin{itemize}
\item the CONUS maximum of $\Tmax$, and the share of western-US land (31--49.4\,$^\circ$N, 124.8--104\,$^\circ$W) with
  $\Tmax\ge40$\,\degC;
\item the area-mean precipitation and the maximum daily-mean wind in four landfall boxes. The boxes are TX coast
  (27--32\,$^\circ$N, 97.8--93.5\,$^\circ$W), FL--GA--SC (26--34\,$^\circ$N, 84--78.5\,$^\circ$W), Louisiana
  (28.8--33\,$^\circ$N, 93.8--88.5\,$^\circ$W) and Southeast (29--37\,$^\circ$N, 86--80\,$^\circ$W);
\item the CONUS maximum and 99th percentile of wind speed.
\end{itemize}
The rule takes \textbf{six event anchors}, each the seasonal maximum of its indicator. To these it adds \textbf{two
fillers} that split the only long event-free gap (5~Aug--11~Sep, 37 days) into thirds. Consecutive dates are therefore
12--16 days apart, close to uniform sampling (92/8 = 11.5 days). The one exception is the deliberate Helene
landfall/inland pair, which lets S3 show the storm moving. Table~\ref{tab:dates} lists the dates and their evidence. All
three SciVis tasks use the same dates (D1--D8) so that results can be compared across tasks.

\figfull{Fig_01_date_selection_timeline.png}{Daily evidence used to choose the eight analysis dates (vertical lines,
labelled D1--D8), gridMET 1~Jul--30~Sep 2024. (a) Hottest CONUS cell (solid) and share of western land with
$\Tmax\ge40$\,\degC{} (dashed). (b) Area-mean daily precipitation in four landfall boxes. (c) Maximum daily-mean 10\,m wind
in the same boxes. (d) CONUS maximum and 99th percentile of daily-mean wind.}{fig:01}

\input{generated/tab_dates}

\textbf{Inferences from Fig.~\ref{fig:01}.}
(i) The heat wave is the dominant July signal. The hottest CONUS cell stays above 52.5\,\degC{} from 4 to 12~July and peaks at
53.85\,\degC{} on 6~July at 36.48\,$^\circ$N, 116.89\,$^\circ$W (Death Valley). The western area above 40\,\degC{} peaks
twice, at about 12--13\,\% around 6--11~July and at 13.6\,\% on 24~July (D2), so D2 is the date with the most
\emph{widespread} heat, not the hottest single cell.
(ii) Each hurricane box has exactly one dominant rain spike, and it matches the landfall date: Texas 40.0\,mm on 8~July (Beryl),
FL--GA--SC 60.1\,mm on 5~August (Debby), Louisiana 54.4\,mm on 11~September (Francine) and the Southeast 93.5\,mm on
26~September (Helene). Helene's area-mean rain is 1.6 times Debby's and covers the largest box. This matches its reported
catastrophic inland rainfall~\cite{nhc_tcr}.
(iii) Daily averaging damps hurricane winds strongly. The highest daily-mean wind of the whole season (16.1\,m/s) occurs
both on 26~September in Florida (Helene) and on 18~September over Wyoming (42.69\,$^\circ$N, 106.06\,$^\circ$W, a
non-tropical mountain-wind event). Only Helene lifts the CONUS 99th percentile above 10\,m/s. Hurricane winds in gridMET
are therefore an order of magnitude below reported peak gusts; we interpret the S3 wind fields as storm-scale circulation,
not as damaging-wind maps.
(iv) Late September is clearly cooler. The CONUS maximum stays at or below 42.6\,\degC{} from 12 to 22~September (minimum 38.2\,\degC{} on 16~September) and the western area
$\ge40$\,\degC{} drops to 0\,\%, so the season moves from heat-dominated in July to storm-dominated in September. This
contrast drives the global-vs-local normalisation discussion in S1.
(v) Day boundaries matter. Helene made landfall as a Category~4 hurricane about 10~miles west-southwest of Perry,
Florida, at 03:10~UTC on 27~September (11:10~pm EDT on the 26th)~\cite{nhc_tcr},
which falls inside gridMET day 26~September. That is why D7 has the landfall rain and wind maximum while D8 shows the inland
decay.

\subsection{The A1 dataset}
The A1 table merges 50 CSVs (871--1,000 posts each) with 19 columns. A1 established, and we re-checked, several facts.
There are no duplicate ids. Timestamps use two formats (48,279 ISO, 987 \texttt{M/D/YYYY H:MM} in AskReddit). Coverage
is strongly skewed in time: 2011--2015 contribute only 778 posts, while 2020 alone contributes 11,493, because the
corpus is the all-time top tail and recent years have had less time to accumulate votes. Five columns are uninformative
or almost empty and are dropped (Sec.~\ref{sec:a1clean}).
